\zadanie{Jądro azotu w polu magnetycznym}

Jądro atomu azotu jest obdarzone momentem magnetycznym, który może przyjmować
trzy orientacje w polu magnetycznym $B$. Stanom tym odpowiadają trzy poziomy
energetyczne o energiach, kolejno,
\begin{equation*}
  \varepsilon_0 = -\gamma_N\hbar B, \qquad
  \varepsilon_1 = 0, \qquad
  \varepsilon_2 = +\gamma_N\hbar B .
\end{equation*}

\begin{enumerate}[label=(\alph*)]
  \item Wyprowadź funkcję podziału $Z(T)$ pojedynczego jądra azotu w polu
    magnetycznym.
  \item Wyprowadź średnią energię $E(T)$ pojedynczego jądra.
  \item Oblicz granice $\lim_{T\to 0} E(T)$ oraz $\lim_{T\to\infty} E(T)$.
  \item Wyprowadź wyrażenia na populacje (obsadzenia) poziomów wzbudzonych
    względem stanu podstawowego, zdefiniowane jako $f_1 = N_1/N_0$
    i $f_2 = N_2/N_0$, gdzie $N_0$ oznacza liczbę jąder w stanie podstawowym $\varepsilon_0$,
    $N_1$ jest liczbą jąder w pierwszym stanie wzbudzonym $\varepsilon_1$,
    a $N_2$ jest liczbą jąder w drugim stanie wzbudzonym $\varepsilon_2$.
  \item Oblicz $f_1$ i $f_2$ w temperaturach:
    (i) $T = 1.0$~K, (ii) $T = 298$~K, (iii) $T = \infty$.
\end{enumerate}

Stałe: $\gamma_N\hbar = 2.04\times 10^{-27}$~J/T, \quad $B = 20$~T, \quad
$k_B = 1.381\times 10^{-23}$~J/K.
